\scanpage{58}
% archive sheet number: 57
\mspage{2 - 3 -}\hfill n\textsuperscript{o} 392

Montrons que c) implique d). Si c) est vraie, l'application $\underset{=}{a}\longrightarrow\operatorname{div}\underset{=}{a}$ est une bijection de $I(A)$ sur $D(A)$, (cf. § 1, n\textsuperscript{o} 1) ; comme c'est un homomorphisme, (§ 1, n\textsuperscript{o} 2), et que $D(A)$ est un groupe, tout élément de $I(A)$ est bien inversible.

Montrons que d) implique e). Si d) est vraie, tout idéal entier $\handwritten{\ne}(0)$ de $A$ est de type fini, (chap. II, § 5, n\textsuperscript{o} 6, th. 4), donc $A$ est noethérien ; comme $I(A)$ est un groupe, $D(A)$ est un groupe, et $A$ est donc complètement intégralement clos, (§ 1, n\textsuperscript{o} 2, th. 1). Enfin, si $\unsure{\underset{=}{p}|\underline{p}|p}$ est un idéal premier non nul de $A$, et si $\underset{=}{m}$ est un idéal maximal de $A$ contenant $\underset{=}{p}$, l'anneau $A_{\underset{=}{m}}$ est principal, (chap. II, § 5, n\textsuperscript{o} 6, th. 4) ; comme $\underset{=}{p} A_{\underset{=}{m}}$ est premier non nul, on a nécessairement $\underset{=}{p} A_{\underset{=}{m}}=\underset{=}{m}A_{\underset{=}{m}}$, d'où $\underset{=}{p}=\underset{=}{m}$, (chap. II, § 2, n\textsuperscript{o} 5, prop. 11), et $\underset{=}{p}$ est maximal.

Montrons que e) implique f). En effet, si $\underset{=}{m}$ est un idéal maximal de $A$, et si e) est vraie, $A_{\underset{=}{m}}$ est un anneau noethérien intégralement clos, et son idéal maximal $\underset{=}{m}A_{\underset{=}{m}}$ est, ou bien $(0)$, ou bien le seul idéal premier non nul de $A_{\underset{=}{m}}$ ; donc $A_{\underset{=}{m}}$ est un corps ou un anneau de valuation discrète d'après la prop. 11 du § 1, n\textsuperscript{o} 7.

Le fait que f) implique g) est évident.

Montrons enfin que g) implique a). Comme $A$ est l'intersection des $A_{\underset{=}{m}}$, ou $\underset{=}{m}$ parcourt l'ensemble des idéaux maximaux, (chap. II, § 3, n\textsuperscript{o} 3, cor. 4 du th. 1), g) implique que $A$ est intégralement clos et noethérien, donc que $A$ est un anneau de Krull, (§ 1, n\textsuperscript{o} 3, cor. du th. 2). D'autre part, on montre que tout idéal premier non nul de $A$ est maximal comme dans la démonstration de d) $\Rightarrow$ e). C. Q. F. D.

PROPOSITION 1.- \uline{Un anneau de Dedekind semi-local est principal}.

Soient $A$ un anneau de Dedekind semi-local, $K$ son corps des fractions, $p_1,\ldots,p_n$ ses idéaux maximaux, et $v_1,\ldots,v_n$ les valuations essentielles correspondantes ; ce sont les seules valuations essentielles de $A$. Soit $\underset{=}{a}$ un idéal entier non nul de $A$. Puisqu'il est divisoriel, il existe, (§ 1, n\textsuperscript{o} 4, prop. 5), des entiers $q_1,\ldots,q_n$ tels que $\underset{=}{a}$ soit l'ensemble des $x\handwritten{\in}K$ tels que $v_i(x)\handwritten{\geq}q_i$ pour
